\documentclass{article}

% ICLR 2025 style approximation
\usepackage[margin=1in]{geometry}
\usepackage{graphicx}
\usepackage{booktabs}
\usepackage{amsmath,amssymb}
\usepackage{enumitem}
\usepackage{hyperref}
\usepackage[table]{xcolor}
\usepackage[font=small]{caption}
\usepackage{subcaption}
\usepackage{natbib}
\usepackage{microtype}
\usepackage{multirow}

\definecolor{grokgreen}{HTML}{2E8B57}
\definecolor{nogrokred}{HTML}{CC3333}
\definecolor{rescueblue}{HTML}{2196F3}

\hypersetup{colorlinks=true, linkcolor=blue, citecolor=blue, urlcolor=blue}

\setlist[itemize]{topsep=2pt,itemsep=1pt,parsep=0pt}
\setlist[enumerate]{topsep=2pt,itemsep=1pt,parsep=0pt}

\title{Grokking Survives Federated Averaging: \\
Phase Boundaries, Aggregation Rescue, and Mechanistic Analysis}

\author{
  Author A \quad Author B \quad Author C \\
  MPhil in Machine Learning, Speech and Language Technology \\
  University of Cambridge \\
  \texttt{\{crsid1, crsid2, crsid3\}@cam.ac.uk}
}

\date{}

\begin{document}
\maketitle

% ============================================================
% ABSTRACT
% ============================================================
\begin{abstract}
Grokking---the phenomenon where neural networks generalise long after memorising the training set---has been extensively studied in centralised settings.
We ask whether grokking survives \emph{federated learning}, where data is distributed across clients and only model updates are aggregated.
Reproducing the setup of \citet{gromov2023grokking} (two-layer MLP with quadratic activations on modular addition), we first establish the centralised grokking phase boundary at a critical training fraction $\alpha_c \approx 0.198$, with grokking time following a power-law divergence $T_{\mathrm{grok}} \propto (\alpha - \alpha_c)^{-0.99}$.
We then conduct over 700 federated training runs across a systematic grid of client counts ($K \in \{2, \ldots, 97\}$), data heterogeneity levels (Dirichlet and structured partitions), and optimisation parameters.
Our main finding is that \textbf{FedAvg aggregation rescues grokking}: in 23 of 24 above-boundary configurations, each client holds data far below the centralised grokking threshold, yet the federated model groks---while a centralised model on the same per-client data never does.
The overhead is modest ($<1.2\times$ slowdown for $K \leq 20$), and heterogeneity delays but does not prevent grokking.
Mechanistic analysis reveals that (i) Fourier feature crystallisation precedes client drift reduction in 100\% of runs, (ii) accuracy, IPR, and drift collapse onto universal curves under time rescaling, and (iii) federation disrupts generalisation timing but not memorisation, identifying a two-timescale separation.
These results demonstrate that grokking is a property of the loss landscape geometry rather than the optimisation path.
\end{abstract}


% ============================================================
% 1. INTRODUCTION
% ============================================================
\section{Introduction}
\label{sec:intro}

\emph{Grokking} refers to the puzzling phenomenon where neural networks first memorise a training set, then---after prolonged additional training---suddenly generalise to unseen data \citep{power2022grokking}.
The effect is most cleanly demonstrated on algorithmic tasks such as modular arithmetic, where \citet{gromov2023grokking} showed that a two-layer MLP with quadratic activation learns periodic Fourier features that explain the generalisation jump.
The emergence of these features is governed by a phase transition: below a critical training fraction $\alpha_c$, the network memorises but never generalises; above it, grokking occurs after a delay that diverges as $\alpha \to \alpha_c^+$.

Meanwhile, \emph{federated learning} (FL) has become the dominant paradigm for privacy-preserving distributed training \citep{mcmahan2017communication, kairouz2021advances}.
In FL, data remains on clients; only model updates are aggregated by a central server.
This introduces several perturbations absent from centralised training: data fragmentation across clients, statistical heterogeneity in local distributions, client drift from multiple local gradient steps, and partial participation.
Each of these could plausibly disrupt the delicate phase transition underlying grokking.

We investigate a simple but fundamental question: \textbf{does grokking survive federated averaging?}
If it does, this implies that the grokking mechanism is robust to substantial perturbations in the optimisation process.
If it does not, the failure mode reveals which aspects of the loss landscape or optimisation dynamics are essential.

Our contributions: (1)~We establish the centralised phase boundary at $\alpha_c \approx 0.198$ with power-law scaling $T_{\mathrm{grok}} \propto (\alpha - \alpha_c)^{-0.99}$ (\S\ref{sec:boundary}).
(2)~In 270 controlled runs, FedAvg groks where per-client data is individually below threshold---aggregation is the sole rescue mechanism (\S\ref{sec:aggregation}).
(3)~Across 147 runs with Dirichlet and structured partitions, heterogeneity delays grokking by up to 37\% near the boundary but does not shift $\alpha_c$ (\S\ref{sec:heterogeneity}).
(4)~Mechanistic analysis of 237 runs reveals that Fourier features crystallise before clients align, observables collapse onto universal curves, and FL selectively disrupts generalisation while preserving memorisation (\S\ref{sec:mechanistic}).


% ============================================================
% 2. BACKGROUND
% ============================================================
\section{Background and Related Work}
\label{sec:background}

\paragraph{Grokking.}
\citet{power2022grokking} first demonstrated grokking on modular arithmetic.
\citet{gromov2023grokking} showed that a minimal 2-layer MLP with quadratic activation suffices, driven by the emergence of periodic Fourier modes, and characterised grokking as a phase transition controlled by the training fraction $\alpha$.
Weight decay \citep{liu2022omnigrok} and representation structure \citep{nanda2023progress} have been identified as key mechanisms; the connection to phase transitions was formalised by \citet{thilak2022slingshot}.

\paragraph{Federated learning.}
FedAvg \citep{mcmahan2017communication} trains a shared model by averaging locally-updated models.
Key challenges include client drift \citep{karimireddy2020scaffold}, non-IID data \citep{li2020federated}, and partial participation; FedProx \citep{li2020fedprox} adds a proximal regulariser to mitigate drift.
To our knowledge, no prior work has studied grokking under federated learning.


% ============================================================
% 3. EXPERIMENTAL SETUP
% ============================================================
\section{Experimental Setup}
\label{sec:setup}

\paragraph{Task and model.}
We study modular addition $(a + b) \bmod p$ with $p = 97$, following \citet{gromov2023grokking}.
The input is a one-hot encoding of $(a, b)$ (dimension $2p = 194$); the output is a distribution over $p$ classes.
We use GrokNet: a two-layer MLP with no biases and quadratic activation $\phi(h) = h^2$, with mean-field initialisation ($W_1 \sim \mathcal{N}(0, 1/D)$, $W_2 \sim \mathcal{N}(0, 1/N^2)$).
The loss is MSE on one-hot targets.
Hidden width is $N = 256$ (validated in preliminary experiments; \S\ref{sec:boundary}).

\paragraph{Training.}
All training uses full-batch gradient descent ($\mathrm{lr} = 50$, no weight decay, no momentum).
Each gradient step processes the entire (local) dataset, so ``epoch'' and ``gradient step'' are synonymous.
We use three seeds (42, 123, 456) and report mean $\pm$ std.

\paragraph{Federated protocol.}
We use Flower \citep{beutel2020flower} with FedAvg.
Default parameters: $K = 5$ clients, $E = 5$ local epochs per round, full participation ($f = 1.0$).
Unless otherwise stated, FL budgets are $S_{\max} = 50{,}000$ gradient steps.

\paragraph{Data partitioning.}
We study four partition strategies:
(i) \textbf{IID}: random equal shards;
(ii) \textbf{Operand}: client $i$ receives samples where $a \bmod K = i$;
(iii) \textbf{Target}: client $i$ receives samples where $(a+b) \bmod K = i$;
(iv) \textbf{Dirichlet}: label-skewed allocation controlled by concentration $\alpha_{\mathrm{dir}}$ (higher $= $ more IID).

\paragraph{Metrics.}
\begin{itemize}
  \item $T_{\mathrm{grok}}$: smallest step where test accuracy $\geq 95\%$ and never drops below (``grokking step'').
  \item $T_{50}$: first step where test accuracy $\geq 50\%$ (onset of generalisation).
  \item $T_{\mathrm{memo}}$: first step where train accuracy $\geq 99\%$ (memorisation).
  \item \textbf{IPR} (Inverse Participation Ratio): measures concentration of Fourier structure in $W_1$; rises from ${\sim}0.02$ (random) to ${\sim}0.35$--$0.48$ during grokking.
  \item \textbf{Client drift}: mean $\|w_{\mathrm{after}} - w_{\mathrm{before}}\|_F$ across clients per round.
\end{itemize}


% ============================================================
% 4. CENTRALISED PHASE BOUNDARY
% ============================================================
\section{Centralised Phase Boundary}
\label{sec:boundary}

We first establish the centralised grokking phase boundary to serve as the reference for all FL experiments.
We validate that hidden width $N = 256$ supports grokking across the $\alpha$ range (all four tested widths $N \in \{100, 128, 200, 256\}$ produce nearly identical $T_{\mathrm{grok}}$; width is not a confound).

\paragraph{Phase boundary.}
Sweeping $\alpha \in \{0.03, 0.05, 0.075, 0.1, 0.15, 0.2, 0.3, 0.5\}$ at $S_{\max} = 100{,}000$ steps reveals a sharp transition (Table~\ref{tab:exp1}): $\alpha \leq 0.2$ shows no grokking (test accuracy $< 7\%$), while $\alpha \geq 0.3$ groks with 3/3 seeds.
Extended runs to $10^6$ steps reveal that $\alpha = 0.2$ \emph{does} eventually grok ($T_{\mathrm{grok}} = 805{,}000$), placing the true critical point below $0.2$.

\begin{table}[t]
\centering
\small
\caption{Centralised phase boundary. $T_{\mathrm{grok}}$ diverges as $\alpha \to \alpha_c^+$.}
\label{tab:exp1}
\begin{tabular}{lcrrc}
\toprule
$\alpha$ & $n_{\mathrm{train}}$ & $T_{\mathrm{grok}}$ (mean $\pm$ std) & Final acc (\%) & Final IPR \\
\midrule
0.10  &    940 & $\infty$ &   0.5 & 0.021 \\
0.15  & 1{,}411 & $\infty$ &   0.6 & 0.022 \\
0.20  & 1{,}881 & $805{,}000^{\dagger}$ & 98.7 & 0.253 \\
0.22  & 2{,}069 & $59{,}000$ & 100.0 & --- \\
0.25  & 2{,}352 & $24{,}000$ & 100.0 & --- \\
0.30  & 2{,}822 & $12{,}833 \pm 404$  & 100.0 & 0.472 \\
0.50  & 4{,}704 & $7{,}667 \pm 115$   & 100.0 & 0.484 \\
\bottomrule
\multicolumn{5}{l}{\footnotesize $\dagger$ Extended to $10^6$ steps; single seed.}
\end{tabular}
\end{table}

\paragraph{Critical scaling.}
Fine-grained runs at $\alpha \in \{0.20, 0.21, 0.22, 0.25, 0.28, 0.30, 0.50\}$ reveal a power-law divergence:
\begin{equation}
  T_{\mathrm{grok}} \propto (\alpha - \alpha_c)^{-\gamma}, \qquad
  \alpha_c \approx 0.198, \quad \gamma \approx 0.99, \quad R^2 = 0.978.
  \label{eq:scaling}
\end{equation}
The exponent $\gamma \approx 1$ means $T_{\mathrm{grok}}$ is roughly inversely proportional to the distance from criticality.
Adding just 94 training samples ($\alpha: 0.20 \to 0.21$) cuts $T_{\mathrm{grok}}$ by $5\times$.
This is consistent with a second-order phase transition: no discontinuity in the order parameter (test accuracy at fixed time), only a diverging timescale.

\paragraph{Mechanistic picture.}
All $\alpha$ values traverse the same dynamical path through (IPR, test accuracy) space---grokking is a universal process; $\alpha$ only controls the clock speed.
The network memorises first (train loss drops, IPR stays flat), then reorganises weights into Fourier structure.
Less data produces a weaker gradient signal for the periodic solution, explaining the critical slowdown.


% ============================================================
% 5. AGGREGATION RESCUES GROKKING
% ============================================================
\section{Aggregation Rescues Grokking}
\label{sec:aggregation}

\paragraph{Experimental design.}
The central question: each FL client holds only $n/K$ samples---often far below $\alpha_c$.
Does FedAvg aggregation recover grokking?
For every $(\alpha, K)$ pair with $\alpha \in \{0.20, 0.25, 0.30, 0.35, 0.50\}$ and $K \in \{2, 5, 10, 20, 50, 97\}$, we run three matched conditions (3 seeds each, 270 runs total):
\begin{enumerate}[nosep]
  \item \textbf{Centralised-full}: trained on all $n_{\mathrm{train}}$ samples (ceiling).
  \item \textbf{Centralised-reduced}: trained on $n_{\mathrm{train}}/K$ samples (floor---what one client could achieve alone).
  \item \textbf{FL IID}: $K$ clients with $n_{\mathrm{train}}/K$ samples each, FedAvg.
\end{enumerate}

\begin{figure}[t]
\centering
\includegraphics[width=\textwidth]{figures/exp2_phase_diagram.png}
\caption{Phase diagram in $(\alpha, K)$ space. \textbf{Gray}: below centralised boundary ($\alpha = 0.2$). \textbf{\textcolor{rescueblue}{Blue}}: FL groks but centralised-reduced does not (aggregation rescues). \textbf{Green}: both grok. \textbf{Red}: neither groks.
Aggregation rescues grokking across the vast majority of the grid.}
\label{fig:exp2_phase}
\end{figure}

\paragraph{Headline result (Figure~\ref{fig:exp2_phase}).}
FedAvg groks in \textbf{23 of 24 above-boundary cells}, even though per-client data is individually below $\alpha_c$.
A centralised model on the same per-client data ($\alpha_{\mathrm{eff}} = \alpha/K$) never groks in any of these cells.
Aggregation is not merely helpful---it is \emph{necessary}: without it, the model never generalises.

\paragraph{Quantifying the overhead (Table~\ref{tab:exp2}).}
The slowdown ratio $T_{\mathrm{grok}}^{\mathrm{FL}} / T_{\mathrm{grok}}^{\mathrm{cent}}$ reveals three regimes: (i) for $K \leq 20$, the overhead is $< 1.2\times$ at all $\alpha$; (ii) at $K = 50$, a moderate $1.01$--$1.17\times$; (iii) at $K = 97$, a significant $1.04$--$2.14\times$, with FL breaking at low $\alpha$.
At $\alpha = 0.5$, FL has essentially zero overhead even at $K = 97$ ($1.04\times$).

\begin{table}[t]
\centering
\small
\caption{FL grokking time and slowdown ratio across the $(\alpha, K)$ grid. Gray: below centralised boundary.}
\label{tab:exp2}
\begin{tabular}{ccrrr}
\toprule
$\alpha$ & $K$ & $T_{\mathrm{grok}}^{\mathrm{cent}}$ & $T_{\mathrm{grok}}^{\mathrm{FL}}$ & Ratio \\
\midrule
\rowcolor{gray!15} 0.20 & 2--97 & $\infty$ & $\infty$ & --- \\
\midrule
0.25 &  2 & $25{,}133$ & $25{,}263$ & $1.01\times$ \\
0.25 & 20 & $25{,}133$ & $29{,}987$ & $1.19\times$ \\
\rowcolor{nogrokred!10} 0.25 & 97 & $25{,}133$ & $\infty$ & --- \\
\midrule
0.30 &  2 & $12{,}833$ & $12{,}822$ & $1.00\times$ \\
0.30 & 50 & $12{,}833$ & $15{,}048$ & $1.17\times$ \\
0.30 & 97 & $12{,}833$ & $27{,}422$ & $2.14\times$ \\
\midrule
0.50 &  2 & $7{,}667$ & $7{,}632$ & $1.00\times$ \\
0.50 & 97 & $7{,}667$ & $7{,}945$ & $1.04\times$ \\
\bottomrule
\end{tabular}
\end{table}

\paragraph{Mechanistic preservation.}
Phase portraits (IPR vs.\ test accuracy) confirm that FL follows the same mechanistic pathway as centralised training: the same S-curve in representation space, the same Fourier solution, at roughly the same speed.
FL does not find a different solution---it finds the same solution via a noisier route.

\paragraph{Client drift dynamics.}
Client drift peaks precisely at the grokking transition, then drops sharply as all clients converge to the Fourier solution.
At the fragmentation limit ($\alpha{=}0.25, K{=}97$), drift grows without bound and the system never finds the attractor.
This drift--grokking interplay is analysed mechanistically in \S\ref{sec:mechanistic}.


% ============================================================
% 6. ROBUSTNESS TO HETEROGENEITY
% ============================================================
\section{Robustness to Heterogeneity}
\label{sec:heterogeneity}

We now test whether data heterogeneity---a central challenge in FL---shifts the grokking phase boundary.

\paragraph{Exp 3a: Dirichlet sweep.}
We sweep $\alpha_{\mathrm{dir}} \in \{0.01, 0.1, 0.5, 1.0, 10.0, 1000.0\}$ (extreme non-IID to IID) at $K = 10$ (Table~\ref{tab:exp3a}).

\begin{table}[t]
\centering
\small
\caption{Dirichlet heterogeneity sweep ($K{=}10$, mean $T_{\mathrm{grok}}$, 3 seeds).}
\label{tab:exp3a}
\begin{tabular}{lcccc}
\toprule
$\alpha$ & $\alpha_{\mathrm{dir}}{=}0.01$ & $0.1$ & $1000$ (IID) & Slowdown \\
\midrule
0.20 & $\infty$ & $\infty$ & $\infty$ & --- \\
0.25 & $36{,}638$ & $29{,}497$ & $26{,}725$ & 1.37$\times$ \\
0.30 & $14{,}448$ & $13{,}690$ & $13{,}072$ & 1.11$\times$ \\
0.50 & $7{,}738$ & $7{,}680$ & $7{,}642$ & 1.01$\times$ \\
\bottomrule
\end{tabular}
\end{table}

The results are clear: \textbf{the phase boundary is unchanged}.
At $\alpha = 0.2$, grokking fails at every heterogeneity level.
At $\alpha \geq 0.25$, all 30 configurations (5 $\alpha_{\mathrm{dir}}$ $\times$ 6 $\alpha$) grok to 100\%.
Heterogeneity delays grokking by up to 37\% at $\alpha = 0.25$ under extreme non-IID ($\alpha_{\mathrm{dir}} = 0.01$), but the effect vanishes in the interior ($<1.3\%$ at $\alpha = 0.5$).

\paragraph{Exp 3b: Structured partitions.}
Operand-based partitioning---which fragments the input space along the Fourier structure the model must learn---was hypothesised to be particularly disruptive.
The data \textbf{does not support this hypothesis}: at $\alpha = 0.5$, all three partitions (IID, operand, target) yield statistically identical $T_{\mathrm{grok}}$ ($7{,}597$--$7{,}642$ steps).
Near the boundary ($\alpha = 0.25$), operand is slightly \emph{faster} than IID ($0.965\times$), while target is slightly slower ($1.114\times$).
The Fourier solution is learned through aggregation regardless of how the input space is divided.

\paragraph{Mechanistic explanation.}
The delay is explained by \textbf{weight growth retardation}: non-IID data slows weight norm growth by ${\sim}20$--$30\%$, delaying Fourier feature emergence.
However, final weight magnitudes and IPR converge to the same values across all heterogeneity levels---drift delays grokking but does not degrade the solution (Appendix~\ref{app:anatomy}).


% ============================================================
% 7. MECHANISTIC ANALYSIS
% ============================================================
\section{Mechanistic Analysis}
\label{sec:mechanistic}

We perform a post-hoc analysis of 237 FL runs from Experiments 2--3 to understand \emph{why} grokking survives federation.
Four findings emerge: (i) FL follows the same mechanistic pathway as centralised training; (ii) client drift predicts grokking delay independently of $\alpha$; (iii) Fourier features crystallise before clients align; and (iv) federation selectively disrupts generalisation while leaving memorisation invariant.

\subsection{Same pathway, different speed}
\label{sec:mech_pathway}

Phase portraits plotting test accuracy against IPR (Figure~\ref{fig:phase_portrait}) reveal that centralised and FL runs trace the \textbf{same L-shaped trajectory} through (IPR, accuracy) space: the model memorises at low IPR, then accuracy jumps once IPR crosses ${\sim}0.12$.
Across all tested $(\alpha, K)$ pairs, FL trajectories are indistinguishable from centralised ones---federation changes the speed along the attractor but not the attractor itself.

Aligning 72 grokking runs to $T_{\mathrm{grok}} = 0$ confirms a \textbf{universal temporal fingerprint} (Appendix~\ref{app:fingerprint}): test accuracy follows a tight sigmoid; the generalisation gap peaks ${\sim}5{,}000$ steps before grokking; client drift follows an inverted-U; and IPR has risen substantially before the accuracy jump.
The narrow interquartile range across runs spanning $K \in \{10, 20\}$, $\alpha_{\mathrm{dir}} \in \{0.01\text{--}1000\}$, and $\alpha \in \{0.25\text{--}0.5\}$ indicates these dynamics are intrinsic to the phase transition.

\begin{figure}[t]
\centering
\includegraphics[width=0.85\textwidth]{figures/exp2_phase_portrait.png}
\caption{Phase portrait: test accuracy vs.\ IPR for centralised and FL IID runs. All trajectories follow the same L-shaped attractor. FL does not alter the pathway---only the speed.}
\label{fig:phase_portrait}
\end{figure}

\subsection{Drift predicts grokking delay, controlling for $\alpha$}
\label{sec:mech_drift}

Naive correlation between drift and $T_{\mathrm{grok}}$ across all runs is $\rho = 0.70$ ($p = 1.6 \times 10^{-29}$), but this is confounded: 42 of 46 failed runs are at $\alpha = 0.2$, which fails centralised too.
Within-$\alpha$ Spearman correlations reveal genuine signal: $\rho = 0.79$ at $\alpha = 0.25$ (strong), declining to $\rho = 0.34$ at $\alpha = 0.5$ (weak).
At $\alpha = 0.3$ with IID partitioning, increasing $K$ from 2 to 97 monotonically increases both drift and $T_{\mathrm{grok}}$ ($\rho = 1.0$), establishing the causal chain $K \to \text{drift} \to T_{\mathrm{grok}}$.

\subsection{Temporal ordering: features crystallise before clients align}
\label{sec:mech_temporal}

Does client consensus enable feature learning, or does feature learning drive consensus?
For each of 72 grokking runs, we detect $T_{\mathrm{IPR\ onset}}$ (IPR exceeds $2.5\times$ baseline) and $T_{\mathrm{drift\ drop}}$ (drift falls below 50\% of peak).
IPR onset precedes drift drop in \textbf{72 of 72 runs} (Figure~\ref{fig:temporal}a; Wilcoxon $p = 1.8 \times 10^{-13}$; median lag 10{,}775 steps), robust across all 12 threshold combinations (Figure~\ref{fig:temporal}d).
The interpretation: the global model discovers periodic Fourier structure first; clients then converge because they are all being pulled toward the same solution.
The Fourier attractor acts as a consensus-inducing fixed point.
Temporal precedence does not establish causation---both could be driven by weight norm growth---but zero exceptions across 72 runs and 12 thresholds constrains alternatives.

\begin{figure}[t]
\centering
\includegraphics[width=\textwidth]{figures/fig8_temporal_ordering.png}
\caption{Temporal ordering of IPR rise and drift drop (72 grokking runs from Exp~3a). (a)~$T_{\mathrm{IPR\ onset}}$ vs.\ $T_{\mathrm{drift\ drop}}$: all points above the diagonal (IPR rises first in 72/72 runs). (b)~Lag histogram (median = 10{,}775 steps; Wilcoxon $p = 1.8 \times 10^{-13}$). (c)~Box plots showing both events precede $T_{\mathrm{grok}}$, with IPR onset leading. (d)~Robustness heatmap: 100\% IPR-first across all 12 threshold combinations.}
\label{fig:temporal}
\end{figure}

\subsection{Scaling collapse}
\label{sec:mech_collapse}

Rescaling time as $\tau = (t - T_{\mathrm{grok}}) / (T_{\mathrm{grok}} - T_{50})$ across 117 grokking runs reveals approximate collapse (Figure~\ref{fig:collapse}).
The accuracy collapse is partly tautological (any sigmoid collapses under centering and scaling), but the IPR and drift collapses are \textbf{non-trivial}: IPR begins rising at $\tau \approx -3$; normalised drift peaks at $\tau \approx -2$.
Since neither was used in the rescaling, their collapse is evidence of a \textbf{single intrinsic timescale} governing the transition.

\begin{figure}[t]
\centering
\includegraphics[width=\textwidth]{figures/fig9_universal_collapse.png}
\caption{Scaling collapse (117 runs). (a)~Raw trajectories. (b)~Rescaled accuracy. (c)~IPR collapse: features emerge at $\tau \approx -3$. (d)~Drift collapse: peaks at $\tau \approx -2$. Black: median; shading: IQR.}
\label{fig:collapse}
\end{figure}

\subsection{Two-timescale separation}
\label{sec:mech_timescales}

Decomposing grokking into memorisation ($T_{\mathrm{memo}}$) and generalisation ($T_{\mathrm{grok}}$) reveals a striking asymmetry (Figure~\ref{fig:two_timescales}):
$T_{\mathrm{memo}}$ is invariant to $K$ ($\pm 1\%$ from $K{=}2$ to $K{=}50$), while $T_{\mathrm{grok}}$ varies by up to $2.1\times$ (CV $4\times$ higher at $\alpha{=}0.25$).
Memorisation requires only that each client overfit its local shard; generalisation requires the global model to discover the Fourier solution via the aggregated gradient signal.
Client drift specifically prolongs the generalisation delay $\Delta T = T_{\mathrm{grok}} - T_{\mathrm{memo}}$ (within-$\alpha$ $\rho \approx 0.73$ at $\alpha{=}0.25$).

\begin{figure}[t]
\centering
\includegraphics[width=\textwidth]{figures/fig10_two_timescales.png}
\caption{FL disrupts generalisation, not memorisation. (a)~Within-$\alpha$ CV: $T_{\mathrm{grok}}$ up to $4\times$ more variable than $T_{\mathrm{memo}}$. (b)~$T_{\mathrm{memo}}$ (dashed, flat) vs.\ $T_{\mathrm{grok}}$ (solid, diverging) vs.\ $K$. (c)~Generalisation delay $\Delta T$ vs.\ mean client drift, with within-$\alpha$ Spearman correlations.}
\label{fig:two_timescales}
\end{figure}


% ============================================================
% 8. DISCUSSION
% ============================================================
\section{Discussion}
\label{sec:discussion}

\paragraph{Grokking is a property of the loss landscape.}
Our central finding is that grokking survives substantial perturbations to the optimisation process: data fragmented across up to 97 clients, non-IID distributions, partial participation, and extreme local epochs.
The mechanistic evidence is multi-layered: identical phase portraits (\S\ref{sec:mech_pathway}), Fourier features driving client consensus (\S\ref{sec:mech_temporal}), and universal scaling collapse (\S\ref{sec:mech_collapse}).

\paragraph{Aggregation as gradient signal amplification.}
Each client's gradient signal for the Fourier solution is individually below threshold, but averaging $K$ signals amplifies the collective signal---analogous to signal averaging in communications.
The two-timescale separation (\S\ref{sec:mech_timescales}) reveals that aggregation does not help memorisation (which occurs locally) but restores the collective signal needed for the Fourier solution.
The Fourier solution itself acts as a \emph{consensus-inducing fixed point}: once periodic features begin crystallising, local gradients align across clients, explaining why drift drops after IPR rises.

\paragraph{Implications for FL practice.}
FedProx's proximal term directly opposes the weight norm growth required for grokking: at $\mu = 0.01$, norms are capped and grokking fails (Appendix~\ref{app:fedprox}).
This reveals a tension between drift-reduction algorithms and beneficial weight dynamics.
Conversely, grokking is robust to extreme local epochs ($E{=}200$ still groks) and partial participation ($f{=}0.2$ caps accuracy at ${\sim}94\%$ but the transition still occurs; Appendix~\ref{app:stress}).

\paragraph{Limitations.}
Restricted to a single architecture (2-layer MLP), activation (quadratic), task (modular addition), and loss (MSE).
We do not save model checkpoints, preventing per-client Fourier spectral analysis.
The temporal ordering establishes precedence, not causation.
All experiments use full-batch training.


% ============================================================
% 10. CONCLUSION
% ============================================================
\section{Conclusion}
\label{sec:conclusion}

We have presented the first systematic study of grokking under federated learning.
Through over 700 training runs spanning client counts, data heterogeneity, optimisation parameters, and FL algorithms, we demonstrate that grokking is remarkably robust to federation.
The grokking phase boundary $\alpha_c \approx 0.198$ is unchanged by federated training, and FedAvg aggregation rescues grokking in configurations where per-client data is individually far below the grokking threshold.

Our mechanistic analysis reveals that grokking under FL follows the same dynamical pathway as centralised training---identical phase portraits, universal temporal fingerprints, and scaling collapse of accuracy, IPR, and drift onto shared curves.
The Fourier solution acts as a consensus-inducing attractor: its crystallisation precedes client alignment in 100\% of runs, and the generalisation timescale is selectively stretched by client drift while memorisation remains invariant.
This multi-layered mechanistic evidence paints a consistent picture: grokking is fundamentally a property of the loss landscape geometry, not the optimisation path.

We hope this work opens several directions: studying grokking under more realistic FL settings (cross-silo with natural data), understanding the aggregation-as-amplification mechanism theoretically, investigating whether per-client Fourier spectra converge to the same modes or merely the same IPR, and exploring whether other phase transitions in learning exhibit similar robustness to distribution.


% ============================================================
% REFERENCES
% ============================================================
\bibliographystyle{plainnat}
\begin{thebibliography}{20}

\bibitem[Beutel et~al.(2020)]{beutel2020flower}
D.~J. Beutel, T.~Tober, D.~Stainton, B.~Roesner, T.~Ollivier, et~al.
\newblock Flower: A friendly federated learning research framework.
\newblock \emph{arXiv preprint arXiv:2007.14390}, 2020.

\bibitem[Gromov(2023)]{gromov2023grokking}
A.~Gromov.
\newblock Grokking modular arithmetic.
\newblock \emph{arXiv preprint arXiv:2301.02679}, 2023.

\bibitem[Kairouz et~al.(2021)]{kairouz2021advances}
P.~Kairouz, H.~B. McMahan, B.~Avent, A.~Bellet, M.~Bennis, et~al.
\newblock Advances and open problems in federated learning.
\newblock \emph{Foundations and Trends in Machine Learning}, 14(1--2):1--210, 2021.

\bibitem[Karimireddy et~al.(2020)]{karimireddy2020scaffold}
S.~P. Karimireddy, S.~Kale, M.~Mohri, S.~J. Reddi, S.~Stich, and A.~T. Suresh.
\newblock {SCAFFOLD}: Stochastic controlled averaging for federated learning.
\newblock In \emph{ICML}, 2020.

\bibitem[Li et~al.(2020{\natexlab{a}})]{li2020federated}
T.~Li, A.~K. Sahu, M.~Zaheer, M.~Sanjabi, A.~Talwalkar, and V.~Smith.
\newblock Federated optimization in heterogeneous networks.
\newblock In \emph{MLSys}, 2020.

\bibitem[Li et~al.(2020{\natexlab{b}})]{li2020fedprox}
T.~Li, A.~K. Sahu, M.~Zaheer, M.~Sanjabi, A.~Talwalkar, and V.~Smith.
\newblock Federated optimization in heterogeneous networks.
\newblock In \emph{MLSys}, 2020.

\bibitem[Liu et~al.(2022)]{liu2022omnigrok}
Z.~Liu, O.~Kitouni, N.~Nolte, E.~J. Michaud, M.~Tegmark, and M.~Williams.
\newblock Omnigrok: Grokking beyond algorithmic data.
\newblock In \emph{ICLR}, 2022.

\bibitem[McMahan et~al.(2017)]{mcmahan2017communication}
B.~McMahan, E.~Moore, D.~Ramage, S.~Hampson, and B.~A. y~Arcas.
\newblock Communication-efficient learning of deep networks from decentralized data.
\newblock In \emph{AISTATS}, 2017.

\bibitem[Nanda et~al.(2023)]{nanda2023progress}
N.~Nanda, L.~Chan, T.~Lieberum, J.~Smith, and J.~Steinhardt.
\newblock Progress measures for grokking via mechanistic interpretability.
\newblock In \emph{ICLR}, 2023.

\bibitem[Power et~al.(2022)]{power2022grokking}
A.~Power, Y.~Burda, H.~Edwards, I.~Babuschkin, and V.~Misra.
\newblock Grokking: Generalization beyond overfitting on small algorithmic datasets.
\newblock In \emph{ICLR (MATH-AI Workshop)}, 2022.

\bibitem[Thilak et~al.(2022)]{thilak2022slingshot}
V.~Thilak, E.~Littwin, S.~Zhai, O.~Saremi, R.~Paiss, and J.~Susskind.
\newblock The slingshot mechanism: An empirical study of adaptive optimizers and the grokking phenomenon.
\newblock \emph{arXiv preprint arXiv:2206.04817}, 2022.

\end{thebibliography}


% ============================================================
% APPENDIX
% ============================================================
\appendix
\section{Full Experiment 2 Results}
\label{app:exp2}

\begin{table}[h]
\centering
\small
\caption{Complete $(\alpha, K)$ grid for Experiment~2. All values are mean over 3 seeds.}
\label{tab:exp2_full}
\begin{tabular}{ccrrrcr}
\toprule
$\alpha$ & $K$ & $\alpha_{\mathrm{eff}}$ & $T_{\mathrm{grok}}^{\mathrm{cent}}$ & $T_{\mathrm{grok}}^{\mathrm{FL}}$ & Grokked & Ratio \\
\midrule
0.25 &  2 & 0.125 & $25{,}133$ & $25{,}263$ & 3/3 & $1.01\times$ \\
0.25 &  5 & 0.050 & $25{,}133$ & $25{,}783$ & 3/3 & $1.03\times$ \\
0.25 & 10 & 0.025 & $25{,}133$ & $26{,}673$ & 3/3 & $1.06\times$ \\
0.25 & 20 & 0.012 & $25{,}133$ & $29{,}987$ & 3/3 & $1.19\times$ \\
0.25 & 50 & 0.005 & $25{,}133$ & --- & 2/3 & --- \\
0.25 & 97 & 0.003 & $25{,}133$ & $\infty$ & 0/3 & --- \\
\midrule
0.30 &  2 & 0.150 & $12{,}833$ & $12{,}822$ & 3/3 & $1.00\times$ \\
0.30 &  5 & 0.060 & $12{,}833$ & $12{,}920$ & 3/3 & $1.01\times$ \\
0.30 & 10 & 0.030 & $12{,}833$ & $13{,}097$ & 3/3 & $1.02\times$ \\
0.30 & 20 & 0.015 & $12{,}833$ & $13{,}480$ & 3/3 & $1.05\times$ \\
0.30 & 50 & 0.006 & $12{,}833$ & $15{,}048$ & 3/3 & $1.17\times$ \\
0.30 & 97 & 0.003 & $12{,}833$ & $27{,}422$ & 3/3 & $2.14\times$ \\
\midrule
0.35 &  2 & 0.175 & $9{,}867$ & $9{,}838$ & 3/3 & $1.00\times$ \\
0.35 &  5 & 0.070 & $9{,}867$ & $9{,}882$ & 3/3 & $1.00\times$ \\
0.35 & 10 & 0.035 & $9{,}867$ & $9{,}952$ & 3/3 & $1.01\times$ \\
0.35 & 20 & 0.017 & $9{,}867$ & $10{,}083$ & 3/3 & $1.02\times$ \\
0.35 & 50 & 0.007 & $9{,}867$ & $10{,}590$ & 3/3 & $1.07\times$ \\
0.35 & 97 & 0.004 & $9{,}867$ & $12{,}837$ & 3/3 & $1.30\times$ \\
\midrule
0.50 &  2 & 0.250 & $7{,}667$ & $7{,}632$ & 3/3 & $1.00\times$ \\
0.50 &  5 & 0.100 & $7{,}667$ & $7{,}640$ & 3/3 & $1.00\times$ \\
0.50 & 10 & 0.050 & $7{,}667$ & $7{,}642$ & 3/3 & $1.00\times$ \\
0.50 & 20 & 0.025 & $7{,}667$ & $7{,}657$ & 3/3 & $1.00\times$ \\
0.50 & 50 & 0.010 & $7{,}667$ & $7{,}737$ & 3/3 & $1.01\times$ \\
0.50 & 97 & 0.005 & $7{,}667$ & $7{,}945$ & 3/3 & $1.04\times$ \\
\bottomrule
\end{tabular}
\end{table}


\section{Centralised Phase Boundary: Extended Results}
\label{app:exp1}

\begin{figure}[h]
\centering
\includegraphics[width=\textwidth]{figures/exp1_mechanistic.png}
\caption{Mechanistic investigation of the centralised phase boundary. (1)~Fourier structure formation rate. (2)~Weight growth drives grokking. (3)~Test loss inflection = grokking onset. (4)~Phase portrait: universal attractor. (5)~Memorise first, then organise. (6)~Power-law fit: $T_{\mathrm{grok}} \propto (\alpha - 0.198)^{-0.99}$.}
\label{fig:exp1_mech}
\end{figure}

\begin{figure}[h]
\centering
\includegraphics[width=\textwidth]{figures/exp1_refined_boundary.png}
\caption{Refined phase boundary with dataset sizes (left), $T_{\mathrm{grok}}$ vs.\ training set size on log-log axes showing power-law scaling (centre), and test accuracy curves for all $\alpha$ values (right).}
\label{fig:exp1_boundary}
\end{figure}

\begin{figure}[h]
\centering
\includegraphics[width=\textwidth]{figures/exp1_ipr_evolution.png}
\caption{IPR evolution across $\alpha$ values (left) and phase portrait showing IPR vs.\ test accuracy (right). All $\alpha$ values that grok converge to the same attractor ($\mathrm{IPR} \approx 0.35$--$0.48$).}
\label{fig:exp1_ipr}
\end{figure}


\section{Aggregation Effect: Additional Figures}
\label{app:exp2_extra}

\begin{figure}[h]
\centering
\includegraphics[width=\textwidth]{figures/exp2_training_curves.png}
\caption{Test accuracy for three representative cells. Blue: centralised-full; red dashed: centralised-reduced; green: FL IID. In panel (b), 10 clients with individually hopeless data shards collectively recover grokking through aggregation alone.}
\label{fig:exp2_curves}
\end{figure}

\begin{figure}[h]
\centering
\includegraphics[width=\textwidth]{figures/exp2_slowdown_ratio.png}
\caption{FL slowdown relative to centralised. Overhead is $<1.2\times$ for $K \leq 20$ and only significant at $K = 97$ near the phase boundary ($\alpha = 0.25$--$0.30$).}
\label{fig:exp2_slowdown}
\end{figure}

\begin{figure}[h]
\centering
\includegraphics[width=\textwidth]{figures/exp2_drift_over_time.png}
\caption{Client drift and test accuracy over time. Drift peaks at the grokking transition and drops once the Fourier solution locks in. At $(\alpha{=}0.25, K{=}97)$, drift grows monotonically.}
\label{fig:exp2_drift}
\end{figure}

\begin{figure}[h]
\centering
\includegraphics[width=\textwidth]{figures/exp2_weight_norms.png}
\caption{Weight norm evolution for layers 1 and 2 across conditions. Centralised-reduced (red dashed) plateaus at low norms---insufficient for grokking. FL IID (green) tracks centralised-full (blue), reaching the same magnitude.}
\label{fig:exp2_wnorms}
\end{figure}

\begin{figure}[h]
\centering
\includegraphics[width=\textwidth]{figures/exp2_fl_vs_centralized.png}
\caption{FL vs.\ centralised $T_{\mathrm{grok}}$ scatter. Points near the diagonal indicate minimal FL overhead; deviation grows with $K$ and lower $\alpha$.}
\label{fig:exp2_scatter}
\end{figure}

\begin{figure}[h]
\centering
\includegraphics[width=\textwidth]{figures/exp2_ipr_evolution.png}
\caption{Fourier structure emergence (IPR) across conditions. FL IID tracks centralised-full; centralised-reduced fails to develop Fourier structure.}
\label{fig:exp2_ipr}
\end{figure}


\section{Grokking Anatomy Across Conditions}
\label{app:anatomy}

\begin{figure}[h]
\centering
\includegraphics[width=\textwidth]{figures/exp3_grokking_anatomy.png}
\caption{Grokking anatomy: four conditions from comfortable ($\alpha{=}0.3$, IID) to near-boundary ($\alpha{=}0.25$, $\alpha_{\mathrm{dir}}{=}0.01$). Rows: test accuracy, generalisation gap, client drift, IPR, weight norms. Heterogeneity stretches the transition by up to 71\% but the final outcome is identical.}
\label{fig:exp3_anatomy}
\end{figure}


\section{Phase-Aligned Universal Fingerprint}
\label{app:fingerprint}

\begin{figure}[h]
\centering
\includegraphics[width=0.65\textwidth]{figures/exp3_phase_aligned_fingerprint.png}
\caption{72 grokking runs aligned to $T_{\mathrm{grok}} = 0$. (a)~Test accuracy, (b)~generalisation gap, (c)~client drift, (d)~IPR. Black: median; shading: IQR. All runs follow a tight common trajectory despite spanning $K \in \{10, 20\}$, $\alpha_{\mathrm{dir}} \in \{0.01\text{--}1000\}$, and $\alpha \in \{0.25\text{--}0.5\}$.}
\label{fig:fingerprint}
\end{figure}


\section{Drift as a Mechanistic Predictor}
\label{app:drift}

\begin{figure}[h]
\centering
\includegraphics[width=\textwidth]{figures/fig7_drift_alpha_controlled.png}
\caption{Drift predicts grokking delay, controlling for $\alpha$. (a)~$T_{\mathrm{grok}}$ vs.\ drift across all 237 runs (red box: 42/46 failures at $\alpha{=}0.2$). (b)~Within-$\alpha$ Spearman $\rho$: significant at every $\alpha$, strongest near the boundary. (c)~Causal chain at $\alpha{=}0.3$: $K \to$ drift $\to T_{\mathrm{grok}}$.}
\label{fig:drift_controlled}
\end{figure}

\begin{figure}[h]
\centering
\includegraphics[width=\textwidth]{figures/fig6_representative_comparison.png}
\caption{Representative run comparison. (a)~Test accuracy, (b)~train loss, (c)~IPR, (d)~mean client drift, (e)~client weight divergence, (f)~layer-1 weight norm. Healthy grok ($\alpha{=}0.5, K{=}10$), delayed grok ($\alpha{=}0.25$, non-IID), and two failure modes (data-limited and fragmented).}
\label{fig:representative}
\end{figure}


\section{Heterogeneity: Additional Results}
\label{app:het}

\begin{figure}[h]
\centering
\includegraphics[width=\textwidth]{figures/exp3a_phase_diagram.png}
\caption{Exp 3a: Dirichlet heterogeneity phase diagram ($K{=}10$). Left: $T_{\mathrm{grok}}$ heatmap. Right: grokking success rate. The phase boundary is unchanged across all heterogeneity levels.}
\label{fig:exp3a_phase}
\end{figure}

\begin{figure}[h]
\centering
\includegraphics[width=\textwidth]{figures/exp3b_structured_partitions.png}
\caption{Exp 3b: Structured partition comparison ($K{=}10$). IID, operand-based, and target-based partitions yield statistically indistinguishable $T_{\mathrm{grok}}$ at all $\alpha$ values.}
\label{fig:exp3b_partitions}
\end{figure}

\begin{figure}[h]
\centering
\includegraphics[width=\textwidth]{figures/exp3_weight_norm_evolution.png}
\caption{Weight norm evolution across heterogeneity levels. Non-IID slows growth by ${\sim}20$--$30\%$ but final values converge.}
\label{fig:exp3_wnorms}
\end{figure}

\begin{figure}[h]
\centering
\includegraphics[width=\textwidth]{figures/exp3_ipr_trajectories.png}
\caption{IPR trajectories by heterogeneity level. More heterogeneous data slows IPR growth but the final plateau is identical.}
\label{fig:exp3_ipr}
\end{figure}

\begin{figure}[h]
\centering
\includegraphics[width=\textwidth]{figures/exp3_grokking_sharpness.png}
\caption{Transition sharpness vs.\ heterogeneity. Left: absolute width ($T_{\mathrm{grok}} - T_{50}$). Right: normalised sharpness. Heterogeneity broadens the transition near the boundary but has minimal effect in the interior.}
\label{fig:exp3_sharpness}
\end{figure}

\begin{figure}[h]
\centering
\includegraphics[width=\textwidth]{figures/exp3_drift_ipr_correlation.png}
\caption{Drift--IPR correlation at the grokking transition. Left: client drift vs.\ IPR at $T_{\mathrm{grok}}$. Right: weight divergence vs.\ IPR. IPR at grokking is invariant to drift level---the same Fourier solution is reached regardless of optimisation noise.}
\label{fig:exp3_drift_ipr}
\end{figure}


\section{Optimisation Fragmentation (Experiment 4)}
\label{app:exp4}

\begin{figure}[h]
\centering
\includegraphics[width=\textwidth]{figures/exp4a_drift.png}
\caption{Exp 4a: Drift accumulation $\times$ heterogeneity ($K{=}10$, $S{=}50$k). At $\alpha{=}0.25$ with non-IID data, $E \geq 25$ prevents grokking (red~$\times$). At $\alpha{=}0.5$, grokking is immune to local epochs.}
\label{fig:exp4a}
\end{figure}

\begin{figure}[h]
\centering
\includegraphics[width=\textwidth]{figures/exp4b_participation.png}
\caption{Exp 4b: Partial participation $\times$ heterogeneity ($K{=}10$, $E{=}5$). Participation fraction has minimal effect on $T_{\mathrm{grok}}$ at all $\alpha$ values, with slight delay under non-IID at $\alpha{=}0.25$.}
\label{fig:exp4b}
\end{figure}

\begin{figure}[h]
\centering
\includegraphics[width=\textwidth]{figures/exp4c_compute_comm.png}
\caption{Exp 4c: Compute vs.\ communication tradeoff ($K{=}10$, $R{=}2000$, IID). At low $\alpha$, insufficient local computation ($E{=}1$) prevents grokking (red~$\times$). At $\alpha{=}0.5$, even $E{=}1$ suffices.}
\label{fig:exp4c}
\end{figure}


\section{FedProx Sweep}
\label{app:fedprox}

\begin{table}[h]
\centering
\small
\caption{FedProx sweep ($K{=}5$, IID, $\alpha{=}0.5$). Sharp transition between $\mu{=}0.001$ (groks) and $\mu{=}0.01$ (fails).}
\label{tab:fedprox}
\begin{tabular}{rlrrrl}
\toprule
$\mu$ & Label & Final test acc & $T_{\mathrm{grok}}$ & $\|W_1\|_F$ & Status \\
\midrule
0     & FedAvg    & 100.0\% & 8{,}215  & 23.80 & Groks \\
0.001 & FedProx   &  99.7\% & 9{,}075  & 22.39 & Groks (10\% delayed) \\
0.01  & FedProx   &   0.1\% & ---   & 15.14 & Weight growth suppressed \\
0.1   & FedProx   &   1.0\% & ---   & NaN   & Diverges \\
\bottomrule
\end{tabular}
\end{table}


\section{Stress Tests: Client Scaling and Extreme Local Epochs}
\label{app:stress}

\begin{table}[h]
\centering
\small
\caption{Client scaling ($\alpha{=}0.5$, IID, full participation). Grokking is immune to client count.}
\begin{tabular}{rrrrr}
\toprule
$K$ & Samples/client & $T_{\mathrm{grok}}$ & Final acc & Final IPR \\
\midrule
5   & 940 & 8{,}215 & 100.0\% & 0.342 \\
10  & 470 & 8{,}220 & 100.0\% & 0.342 \\
20  & 235 & 8{,}240 & 100.0\% & 0.340 \\
50  & 94  & 8{,}280 & 100.0\% & 0.332 \\
97  & 48  & 8{,}430 & 100.0\% & 0.311 \\
\bottomrule
\end{tabular}
\end{table}

\begin{table}[h]
\centering
\small
\caption{Extreme local epochs ($K{=}5$, IID). Even $E{=}200$ (100 rounds) groks. Higher $E$ gives higher IPR.}
\begin{tabular}{rrrrr}
\toprule
Local epochs ($E$) & Rounds & $T_{\mathrm{grok}}$ & Final acc & Final IPR \\
\midrule
50  & 400 & 8{,}300  & 100.0\% & 0.457 \\
100 & 200 & 8{,}400  & 100.0\% & 0.446 \\
200 & 100 & 9{,}000  & 100.0\% & 0.426 \\
\bottomrule
\end{tabular}
\end{table}

\begin{table}[h]
\centering
\small
\caption{Partial participation ($K{=}20$, target partition, $E{=}5$). The only axis that degrades accuracy.}
\begin{tabular}{rrrrrr}
\toprule
Fraction ($f$) & Clients/round & $T_{\mathrm{grok}}$ & Final acc & Final IPR \\
\midrule
1.0 & 20 & 9{,}585 & 98.5\% & 0.214 \\
0.8 & 16 & 9{,}600 & 97.5\% & 0.218 \\
0.6 & 12 & 9{,}725 & 96.2\% & 0.220 \\
0.4 &  8 & 9{,}900 & 93.9\% & 0.222 \\
0.2 &  4 & 9{,}720 & 94.5\% & 0.242 \\
\bottomrule
\end{tabular}
\end{table}


\section{Width Validation (Experiment 0)}
\label{app:exp0}

\begin{figure}[h]
\centering
\includegraphics[width=\textwidth]{figures/exp0_width_validation.png}
\caption{Width validation: $N \in \{100, 128, 200, 256\}$ across $\alpha$ values. All widths produce nearly identical $T_{\mathrm{grok}}$; width is not a confound.}
\label{fig:exp0}
\end{figure}


\section{Seed Consistency}
\label{app:seeds}

\begin{figure}[h]
\centering
\includegraphics[width=\textwidth]{figures/exp2_seed_consistency.png}
\caption{Exp 2: Seed consistency across the $(\alpha, K)$ grid. Low variance at high $\alpha$; increasing variability near the phase boundary.}
\label{fig:exp2_seeds}
\end{figure}


\end{document}
